% TOP_PROBLEM: {"id": 5200006, "title": "Open Problems on Billiards and Geometric Optics", "category_id": 11, "category": {"id": 11, "name": "dynamical_systems", "display_name": "Dynamical Systems", "description": "Problems about long-term behavior of deterministic systems, Hamiltonian dynamics, and periodic orbits.", "slug": "dynamical-systems", "order_index": 11, "created_at": "2026-07-31T15:26:25.671Z"}, "status": "partially_solved"}
% FIRST_POSTED: null
% ENTRY_KIND: statement_only
% Imported from: batch03.tex; UnsolvedMath ID 5200006
\documentclass[11pt]{article}
\usepackage[margin=1in]{geometry}
\usepackage{amsmath,amssymb,mathtools}
\usepackage{fontspec,unicode-math}
\setmainfont{Latin Modern Roman}
\setsansfont{Latin Modern Sans}
\setmathfont{Latin Modern Math}
\usepackage{xurl,hyperref}
\hypersetup{colorlinks=true,urlcolor=blue,linkcolor=blue}
\newcommand{\sref}[2]{\hyperlink{source:#1:#2}{[S#2]}}
\newcommand{\cataloguescope}[1]{\paragraph{Recorded mathematical question.}#1}
\begin{document}
% UnsolvedMath ID 5200006; source catalogue code AMR-051-0006
\setcounter{section}{5200005}
\section{Open Problems on Billiards and Geometric Optics}\label{5200006}
{\small\sffamily UnsolvedMath ID 5200006 · Dynamical Systems}
\subsection{Definitions and mathematical statement}

\paragraph{Shared notation.}$\mathbb N=\{1,2,\ldots\}$, and $\mathbb Z,\mathbb Q,\mathbb R,\mathbb C$ denote the integers, rationals, reals and complex numbers. Any other conventions are specified locally.


The line field $L$ is smooth on the smooth boundary pieces and is transverse to their tangent lines. A harmonic quadruple in the pencil through $p$ means that the nonidentity projective involution fixing $L(p)$ and $T_p\partial\Omega$ exchanges the incident and reflected lines (equivalently their cross-ratio is $-1$ in this ordering). The reflected ray is oriented back into the domain. A $k$-periodic point returns after $k$ applications of the billiard map; $k$-reflective means that such points include a nonempty open subset.
\paragraph{Statement recorded in the supplied source.}
A planar projective billiard is a bounded domain $\Omega$ whose piecewise-smooth boundary carries a transverse line field $L$. At $p\in\partial\Omega$, an incident line $\ell$ reflects to $\ell'$ when $\ell,\ell',L(p),T_p\partial\Omega$ form a harmonic quadruple. Call the billiard $k$-reflective if its billiard map has an open set of $k$-periodic points. (1) Construct a $k$-reflective projective billiard for some odd $k\ge5$. (2) For fixed $k\ge4$, classify $k$-reflective projective billiards within natural boundary-smoothness classes (polygonal, piecewise algebraic, analytic, and so on)—the projective Ivrii conjecture. \sref{5200006}{2}
\subsection{Short English statement}
Find an odd-period reflective projective billiard with period at least five; for each period at least four, determine the examples within specified boundary regularity classes.
\subsection{Sources}
\begin{description}
\item[\hypertarget{source:5200006:1}{S1}] ULAM.AI and the UnsolvedMath Contributors, \emph{UnsolvedMath: A Curated Collection of Open Mathematics Problems}, record 5200006 (AMR-051-0006). CC BY 4.0. \url{https://huggingface.co/datasets/ulamai/UnsolvedMath}.
\item[\hypertarget{source:5200006:2}{S2}] M. Bialy, C. Fierobe, A. Glutsyuk, M. Levi, A. Plakhov and S. Tabachnikov, \emph{Open problems on billiards and geometric optics} (2021), arXiv:2110.10750v2, Corentin Fierobe, Problem 1, pp. 4--5. \url{https://arxiv.org/pdf/2110.10750v2}.
\end{description}
\label{end:5200006}
\end{document}
